\documentclass{article}
\usepackage{geometry}
\geometry{a4paper, margin=1in}
\usepackage{graphicx}
\usepackage{hyperref}

\title{Practical Machine Learning and Deep Learning\\D1.1 Progress Submission\\
Team "Magic Photo"}
\author{Maxim Ilin \href{mailto:m.ilin@innopolis.university}{m.ilin@innopolis.university},\\ Rail Sabirov \href{mailto:rai.sabirov@innopolis.university}{rai.sabirov@innopolis.university},\\ Ivan Ilyichev \href{mailto:i.ilyichev@innopolis.university}{i.ilyichev@innopolis.university}}

\date{\today}

\begin{document}

\maketitle

\section{Project Topic}
\par This project objective is to enhance existing solution for image restoration from different corruptions and coloring of grayscale images using machine learning and deep learning techniques. The project will focus on developing a user-friendly application that allows users to upload images, select the processing option they want to apply, and receive the processed image as output. Additionally, we plan to collect users' feedback for future product improvements and keep track of application lifecycle. The application will leverage pre-trained models and fine-tune them. 
\par Target users are everyone because we want to develop a user-friendly application. The value is that it will allow to restore and colorize not only images, but also important and warm moments of life that were captured in photos.
\par We assume that for coloring people will use mostly faces, cities, pets, and nature images. For restoration this assumption is not so obvious and strict, the exact focus is under discussion now.

\section{Link to GitHub repository}
\par The link to the GitHub repository: \url{https://github.com/ivaniinno/Image-Restoration-and-Coloring}

\section{Existing Solutions}
\par We discovered several promising solutions:
\begin{itemize}
    \item \url{https://www.kaggle.com/code/mdzaidanwar/simple-gfpgan-implementation} - A simple implementation of Generative Facial Prior Generative Adversarial Network (GFPGAN) for image restoration.
    \item 
\end{itemize}

\section{Dataset}
\par Since we focus on people, we would like to use the following dataset with photo links and captions: \url{https://huggingface.co/datasets/nickpai/coco2017-colorization/viewer/default/train?row=18&views%5B%5D=train} (and a version with photos only: \url{https://huggingface.co/datasets/ayushshah/coco-2017-image-colorization-224}). 
\par Also, a sample (for performance reason) of \textbf{celebfaces} is a good option.

\section{Project Succes Criteria}
\par We want to compare the results of models right out of the box and after fine-tuning from our side. The good result would be to achieve at least \(+10\%\) in accuracy and to have at least \(6/10\) positive feedbacks from users on average.

\section{Work Distibution}
\begin{itemize}
    \item Maxim Ilin: finding and processing datasets
    \item Rail Sabirov: MLOps, API, and Frontend
    \item Ivan Ilyichev: GitHub maintenance, reports, slides
    \item Everyone: finding and fine-tuning appropriate models with proper logging
\end{itemize}

\section{Project Sharing among Courses}
\par For now we do not plan to share this project with other courses.


\end{document}